% GENERATED FILE. Edit canonical lessons, then run npm run generate:books.
% canonical-source-hash: fnv1a64:4d598fdf8b8380cc
% canonical-lessons: JA-C340-jogen, JA-C340-chuukoku-suru, JA-C340-kokorogakeru, JA-C340-doryoku-suru, JA-C340-keikoku, JA-R340-review-335-340

\chapter{Giving Advice}
\label{ch:ja-giving-advice}
\hlchaptermodality{340}

\begin{chapteropening}
\textbf{By the end of this chapter:} \emph{I can give advice and say why it is good advice.}

Complete the last lesson of chapter 340: I can give advice and say why it is good advice.
\end{chapteropening}

\section[\texorpdfstring{\ja{じょげん} (jogen)}{じょげん (jogen)}]{\ja{じょげん} (jogen) — advice}

\label{lesson:JA-C340-jogen}



\begin{quote}
Before the new one: say the Japanese for to adopt an idea, then the Japanese for scattered.
\end{quote}

\subsection*{You'll want to know: \ja{じょげん}}
\textbf{\ja{じょげん}} — \emph{jogen} — \textquotedblleft{}advice\textquotedblright{}.

\begin{quote}
\textbf{\ja{せんぱいの} \ja{じょげんを} \ja{きいた} \ja{ほうが} \ja{いいです。}} — \emph{senpai no jogen o kiita hō ga ii desu} — \textquotedblleft{}You should listen to your senior's advice.\textquotedblright{}
\end{quote}

\begin{cousinweb}
\textbf{\ja{じょげん}} (\emph{jogen}) is Sino-Japanese (kango): two Chinese characters read with their on'yomi, readings borrowed from Chinese. They are \emph{jo}, \textquotedblleft{}help\textquotedblright{} (Middle Chinese \emph{dzrjoH}) and \emph{gen}, \textquotedblleft{}word\textquotedblright{} (\emph{ngjon}) (KANJIDIC2; Middle Chinese in Baxter's transcription).
\end{cousinweb}

\begin{grammarlens}[title={What you've built}]
One more word for giving reasons and opinions.
\end{grammarlens}

\subsection*{Your turn}
\begin{itemize}
\raggedright
  \item \emph{Say it:} \emph{jogen}
  \item \emph{Say it:} \emph{jogen}, once more
  \item \emph{Say it:} \emph{senpai no jogen o kiita hō ga ii desu} — You should listen to your senior's advice.
  \item \emph{Recall:} say the Japanese for to adopt an idea, then the Japanese for scattered, then say \emph{jogen} again
\end{itemize}

\begin{tcolorbox}[breakable,colback=teal!4,colframe=teal!35!black,title={Before you move on}]
What does \ja{じょげん} mean? (\textquotedblleft{}Advice\textquotedblright{}.) Say it once more.
\end{tcolorbox}
\section[\texorpdfstring{\ja{ちゅうこくする} (chūkoku suru)}{ちゅうこくする (chūkoku suru)}]{\ja{ちゅうこくする} (chūkoku suru) — to warn; to advise against}

\label{lesson:JA-C340-chuukoku-suru}



\begin{quote}
Before the new one: say the Japanese for scattered, then the Japanese for advice.
\end{quote}

\subsection*{You'll want to know: \ja{ちゅうこくする}}
\textbf{\ja{ちゅうこくする}} — \emph{chūkoku suru} — \textquotedblleft{}to warn; to advise against\textquotedblright{}.

\begin{quote}
\textbf{\ja{あぶないから、いしゃが} \ja{ちゅうこくしました。}} — \emph{abunai kara, isha ga chūkoku shimashita} — \textquotedblleft{}It is dangerous, so the doctor warned me.\textquotedblright{}
\end{quote}

Polite \textbf{\ja{ちゅうこくします}}, \emph{chūkokushimasu}; a \ja{する} verb. Past (polite): \textbf{\ja{いしゃが} \ja{ちゅうこくしました。}}, \emph{isha ga chūkoku shimashita}.

\begin{cousinweb}
\textbf{\ja{ちゅうこくする}} (\emph{chūkoku suru}) is built from a Sino-Japanese noun and the native verb \emph{suru}, \textquotedblleft{}to do\textquotedblright{}. The noun, \emph{chūkoku}, is two Chinese characters read with their on'yomi, readings borrowed from Chinese: \emph{chū}, \textquotedblleft{}loyalty\textquotedblright{} and \emph{koku}, \textquotedblleft{}tell\textquotedblright{} (Middle Chinese \emph{kowk}) (KANJIDIC2; JMdict; Middle Chinese in Baxter's transcription).
\end{cousinweb}

\begin{grammarlens}[title={What you've built}]
One more word for giving reasons and opinions.
\end{grammarlens}

\subsection*{Your turn}
\begin{itemize}
\raggedright
  \item \emph{Say it:} \emph{chūkoku suru}
  \item \emph{Say it:} \emph{chūkoku suru}, once more
  \item \emph{Say it:} \emph{abunai kara, isha ga chūkoku shimashita} — It is dangerous, so the doctor warned me.
  \item \emph{Recall:} say the Japanese for scattered, then the Japanese for advice, then say \emph{chūkoku suru} again
\end{itemize}

\begin{tcolorbox}[breakable,colback=teal!4,colframe=teal!35!black,title={Before you move on}]
What does \ja{ちゅうこくする} mean? (\textquotedblleft{}To warn; to advise against\textquotedblright{}.) Say it once more.
\end{tcolorbox}
\section[\texorpdfstring{\ja{こころがける} (kokorogakeru)}{こころがける (kokorogakeru)}]{\ja{こころがける} (kokorogakeru) — to keep in mind; to make a point of}

\label{lesson:JA-C340-kokorogakeru}



\begin{quote}
Before the new one: say the Japanese for advice, then the Japanese for to warn.
\end{quote}

\subsection*{You'll want to know: \ja{こころがける}}
\textbf{\ja{こころがける}} — \emph{kokorogakeru} — \textquotedblleft{}to keep in mind; to make a point of\textquotedblright{}.

\begin{quote}
\textbf{\ja{けんこうの} \ja{ために、はやく} \ja{ねる} \ja{ことを} \ja{こころがけて} \ja{ください。}} — \emph{kenkō no tame ni, hayaku neru koto o kokorogakete kudasai} — \textquotedblleft{}For your health, make a point of going to bed early.\textquotedblright{}
\end{quote}

Polite \textbf{\ja{こころがけます}}, \emph{kokorogakemasu}; an ichidan verb. Past (polite): \textbf{\ja{あには} \ja{あさの} \ja{うんどうを} \ja{こころがけました。}}, \emph{ani wa asa no undō o kokorogakemashita}.

\begin{cousinweb}
\textbf{\ja{こころがける}} (\emph{kokorogakeru}) is a native Japanese word (wago), inherited rather than borrowed from Chinese: it is built from \emph{kokoro}, \textquotedblleft{}heart\textquotedblright{} and \emph{kakeru}, \textquotedblleft{}to hang\textquotedblright{}, and its kanji were attached later as kun'yomi (JMdict).
\end{cousinweb}

\begin{grammarlens}[title={What you've built}]
One more word for giving reasons and opinions.
\end{grammarlens}

\subsection*{Your turn}
\begin{itemize}
\raggedright
  \item \emph{Say it:} \emph{kokorogakeru}
  \item \emph{Say it:} \emph{kokorogakeru}, once more
  \item \emph{Say it:} \emph{kenkō no tame ni, hayaku neru koto o kokorogakete kudasai} — For your health, make a point of going to bed early.
  \item \emph{Recall:} say the Japanese for advice, then the Japanese for to warn, then say \emph{kokorogakeru} again
\end{itemize}

\begin{tcolorbox}[breakable,colback=teal!4,colframe=teal!35!black,title={Before you move on}]
What does \ja{こころがける} mean? (\textquotedblleft{}To keep in mind; to make a point of\textquotedblright{}.) Say it once more.
\end{tcolorbox}
\section[\texorpdfstring{\ja{どりょくする} (doryoku suru)}{どりょくする (doryoku suru)}]{\ja{どりょくする} (doryoku suru) — to make an effort}

\label{lesson:JA-C340-doryoku-suru}



\begin{quote}
Before the new one: say the Japanese for to warn, then the Japanese for to make a point of.
\end{quote}

\subsection*{You'll want to know: \ja{どりょくする}}
\textbf{\ja{どりょくする}} — \emph{doryoku suru} — \textquotedblleft{}to make an effort\textquotedblright{}.

\begin{quote}
\textbf{\ja{しけんが} \ja{あるから、どりょくした} \ja{ほうが} \ja{いいです。}} — \emph{shiken ga aru kara, doryoku shita hō ga ii desu} — \textquotedblleft{}There is an exam, so you should make an effort.\textquotedblright{}
\end{quote}

Polite \textbf{\ja{どりょくします}}, \emph{doryokushimasu}; a \ja{する} verb. Past (polite): \textbf{\ja{がくせいが} \ja{どりょくしました。}}, \emph{gakusei ga doryoku shimashita}.

\begin{cousinweb}
\textbf{\ja{どりょくする}} (\emph{doryoku suru}) is built from a Sino-Japanese noun and the native verb \emph{suru}, \textquotedblleft{}to do\textquotedblright{}. The noun, \emph{doryoku}, is two Chinese characters read with their on'yomi, readings borrowed from Chinese: \emph{do}, \textquotedblleft{}toil\textquotedblright{} (Middle Chinese \emph{nuX}) and \emph{ryoku}, \textquotedblleft{}power\textquotedblright{} (\emph{lik}) (KANJIDIC2; JMdict; Middle Chinese in Baxter's transcription).
\end{cousinweb}

\begin{grammarlens}[title={What you've built}]
One more word for giving reasons and opinions.
\end{grammarlens}

\subsection*{Your turn}
\begin{itemize}
\raggedright
  \item \emph{Say it:} \emph{doryoku suru}
  \item \emph{Say it:} \emph{doryoku suru}, once more
  \item \emph{Say it:} \emph{shiken ga aru kara, doryoku shita hō ga ii desu} — There is an exam, so you should make an effort.
  \item \emph{Recall:} say the Japanese for to warn, then the Japanese for to make a point of, then say \emph{doryoku suru} again
\end{itemize}

\begin{tcolorbox}[breakable,colback=teal!4,colframe=teal!35!black,title={Before you move on}]
What does \ja{どりょくする} mean? (\textquotedblleft{}To make an effort\textquotedblright{}.) Say it once more.
\end{tcolorbox}
\section[\texorpdfstring{\ja{けいこく} (keikoku)}{けいこく (keikoku)}]{\ja{けいこく} (keikoku) — a warning}

\label{lesson:JA-C340-keikoku}



\begin{quote}
Before the new one: say the Japanese for to make a point of, then the Japanese for to make an effort.
\end{quote}

\subsection*{You'll want to know: \ja{けいこく}}
\textbf{\ja{けいこく}} — \emph{keikoku} — \textquotedblleft{}a warning\textquotedblright{}.

\begin{quote}
\textbf{\ja{けいこくを} \ja{きいたので、みんな} \ja{うちに} \ja{かえりました。}} — \emph{keikoku o kiita node, minna uchi ni kaerimashita} — \textquotedblleft{}Everyone heard the warning, so they went home.\textquotedblright{}
\end{quote}

\begin{cousinweb}
\textbf{\ja{けいこく}} (\emph{keikoku}) is Sino-Japanese (kango): two Chinese characters read with their on'yomi, readings borrowed from Chinese. They are \emph{kei}, \textquotedblleft{}admonish\textquotedblright{} (Middle Chinese \emph{kjaengX}) and \emph{koku}, \textquotedblleft{}tell\textquotedblright{} (\emph{kowk}) (KANJIDIC2; Middle Chinese in Baxter's transcription).
\end{cousinweb}

\begin{grammarlens}[title={What you've built}]
One more word for giving reasons and opinions.
\end{grammarlens}

\subsection*{Your turn}
\begin{itemize}
\raggedright
  \item \emph{Say it:} \emph{keikoku}
  \item \emph{Say it:} \emph{keikoku}, once more
  \item \emph{Say it:} \emph{keikoku o kiita node, minna uchi ni kaerimashita} — Everyone heard the warning, so they went home.
  \item \emph{Recall:} say the Japanese for to make a point of, then the Japanese for to make an effort, then say \emph{keikoku} again
\end{itemize}

\begin{tcolorbox}[breakable,colback=teal!4,colframe=teal!35!black,title={Before you move on}]
What does \ja{けいこく} mean? (\textquotedblleft{}A warning\textquotedblright{}.) Say it once more.
\end{tcolorbox}
\section[(dialogue)]{Looking back: from Groups and Members to Giving Advice}

\label{lesson:JA-R340-review-335-340}



\begin{quote}
Say the Japanese for an organization and a club member.
\end{quote}

\subsection*{Your turn}
Say these aloud:

\begin{itemize}
\raggedright
  \item an organization, a club member, a membership fee, oneself, other people
  \item the environment, nature, the earth, air, pollution
  \item to economize, electric power, fuel, petroleum, too good to waste
  \item useless, to reduce, to increase something, to burn something, resources
  \item to sort rubbish, to preserve, to safeguard, to adopt an idea, scattered
  \item advice, to warn, to make a point of, to make an effort, a warning
\end{itemize}

\begin{tcolorbox}[breakable,colback=teal!4,colframe=teal!35!black,title={Before you move on}]
An organization? (\textbf{\ja{だんたい}}, \emph{dantai}.) A club member? (\textbf{\ja{かいいん}}, \emph{kaiin}.) A membership fee? (\textbf{\ja{かいひ}}, \emph{kaihi}.) Oneself? (\textbf{\ja{じぶん}}, \emph{jibun}.) Other people? (\textbf{\ja{たにん}}, \emph{tanin}.) The environment? (\textbf{\ja{かんきょう}}, \emph{kankyō}.) Nature? (\textbf{\ja{しぜん}}, \emph{shizen}.) The earth? (\textbf{\ja{ちきゅう}}, \emph{chikyū}.) Air? (\textbf{\ja{くうき}}, \emph{kūki}.) Pollution? (\textbf{\ja{おせん}}, \emph{osen}.) To economize? (\textbf{\ja{せつやくする}}, \emph{setsuyaku suru}.) Electric power? (\textbf{\ja{でんりょく}}, \emph{denryoku}.) Fuel? (\textbf{\ja{ねんりょう}}, \emph{nenryō}.) Petroleum? (\textbf{\ja{せきゆ}}, \emph{sekiyu}.) Too good to waste? (\textbf{\ja{もったいない}}, \emph{mottainai}.) Useless? (\textbf{\ja{むだ}}, \emph{muda}.) To reduce? (\textbf{\ja{へらす}}, \emph{herasu}.) To increase something? (\textbf{\ja{ふやす}}, \emph{fuyasu}.) To burn something? (\textbf{\ja{もやす}}, \emph{moyasu}.) Resources? (\textbf{\ja{しげん}}, \emph{shigen}.) To sort rubbish? (\textbf{\ja{ぶんべつする}}, \emph{bunbetsu suru}.) To preserve? (\textbf{\ja{ほぞんする}}, \emph{hozon suru}.) To safeguard? (\textbf{\ja{ほごする}}, \emph{hogo suru}.) To adopt an idea? (\textbf{\ja{とりいれる}}, \emph{toriireru}.) Scattered? (\textbf{\ja{ばらばら}}, \emph{barabara}.) Advice? (\textbf{\ja{じょげん}}, \emph{jogen}.) To warn? (\textbf{\ja{ちゅうこくする}}, \emph{chūkoku suru}.) To make a point of? (\textbf{\ja{こころがける}}, \emph{kokorogakeru}.) To make an effort? (\textbf{\ja{どりょくする}}, \emph{doryoku suru}.) A warning? (\textbf{\ja{けいこく}}, \emph{keikoku}.)
\end{tcolorbox}
